% Multi-copy repeats: where the two-copy theory of assembly ambiguity stops
% Build: pdflatex multicopy (x3)
\documentclass[11pt]{article}
\usepackage{lmodern}
\usepackage[T1]{fontenc}
\usepackage{amsmath,amssymb,amsthm,mathtools}
\usepackage[margin=1.1in]{geometry}
\usepackage{booktabs}
\usepackage{microtype}
\usepackage[hidelinks]{hyperref}

\theoremstyle{plain}
\newtheorem{theorem}{Theorem}
\newtheorem{proposition}[theorem]{Proposition}
\newtheorem{lemma}[theorem]{Lemma}
\newtheorem{corollary}[theorem]{Corollary}
\theoremstyle{definition}
\newtheorem{definition}[theorem]{Definition}
\theoremstyle{remark}
\newtheorem{remark}[theorem]{Remark}

\newcommand{\Z}{\mathbb{Z}}
\newcommand{\K}{\mathcal{K}}
\newcommand{\coker}{\operatorname{coker}}
\newcommand{\Sk}{\mathrm{Sk}}

\title{\bfseries Multi-copy repeats: where the two-copy theory\\ of assembly ambiguity stops}
\author{Ruqing Chen\\[2pt]
\normalsize GUT Geoservice Inc., Rivi\`ere-Beaudette, Qu\'ebec J0P 1R0, Canada\\
\normalsize\texttt{ruqing@hotmail.com}}
\date{25 September 2026}

\begin{document}
\maketitle

\begin{abstract}
\noindent
The sandpile group of a genome's de Bruijn graph splits into a length-driven part and
the group $\K(D_w)$ of the transition digraph of the circular word of repeat occurrences
\cite{Bio3}. For repeats of multiplicity two, $\K(D_w)$ is the critical group of a
one-vertex map: $\coker(A+I)$ for Bouchet's signed interlace matrix, by a theorem of
Merino, Moffatt and Noble --- a presentation whose entries are \emph{pairwise} data
about the repeats. This note asks what survives for repeats with three or more copies.
Two families of any multiplicities $r,s$ give $\K(D_w)=\Z/c$, where $c$ is the number of
transitions from one family to the other, $1\le c\le\min(r,s)$; this contains the
interleaving theorem for two-copy repeats. Beyond two families nothing pairwise survives:
for three-copy repeats, two words with identical pairwise occurrence patterns already
have groups $\Z/3$ and $\Z/4$ at $m=3$, so no matrix with pairwise entries can present
$\K(D_w)$; detaching a three-copy vertex into two-copy vertices does not control the
group; and the literature has no interlacement formula for $r$-in $r$-out digraphs with
$r\ge3$. A census of the groups realised by three- and four-copy arrangements is given
($9$, $17$ and $28$ groups for $(3,3,3)$, $(4,4,4)$ and $(3,3,3,3)$, up to $(\Z/3)^{3}$).
On $\varphi$X174 the skeleton is $2$-in $2$-out at $k=11,10$, where the signed interlace
matrix reproduces the group exactly, and acquires vertices of out-degree $3$ to $6$ from
$k=9$ down: the two-copy presentation stops exactly where the arrangement group becomes
large. Four checks, all exact.
\end{abstract}

\section{Introduction}\label{sec:intro}

Let $S$ be a circular sequence with repeat families $W_1,\dots,W_m$ of multiplicities
$r_i$, clean in the sense of \cite{Bio1}, and let $w$ be the circular word of their
occurrences. The bundle-chain splitting theorem \cite{Bio3} gives
\[
\K\bigl(G_k(S)\bigr)\;\cong\;\bigoplus_i(\Z/r_i)^{\,s_i-1}\;\oplus\;\K(D_w),
\]
where $D_w$ has one vertex per family and one arc per adjacency of occurrences. The
first summand is what longer $k$-mers remove. The second, the \emph{arrangement group},
is what no $k$ removes, and for $r_i=2$ it has a striking description: $D_w$ is the
directed medial graph of a bouquet whose edges are the repeats, and by
\cite[Thm.~7.3]{MMN} $\K(D_w)\cong\coker(A+I)$ with $A$ the skew-symmetric
$\{0,\pm1\}$ matrix recording, for each pair of repeats, whether and in which order
their occurrences interlace. The entries of $A$ are pairwise data; its determinant is the
interlace polynomial at $1$ \cite{ABS}, the number of Euler circuits of $D_w$
\cite{MacrisPule,Jonsson}.

Real genomes have repeats with more than two copies. In $\varphi$X174 the skeleton of
\cite{Bio3} is $2$-in $2$-out at $k=11$ and $k=10$ but has vertices of out-degree $3$ and
$4$ at $k=9$ and up to $6$ at $k=7$ (Table~\ref{tab:phix}). This note asks whether the
two-copy description extends. The answer is: only in the smallest case.

\paragraph{What is known.} For $2$-in $2$-out digraphs the determinant formula
$\#\text{Euler circuits}=\det(I+A)$ is due to Macris and Pul\'e \cite{MacrisPule}, with
elementary proofs by Lauri and Jonsson \cite{Jonsson}; the group statement is
\cite{MMN}. Jonsson \cite[p.~192]{Jonsson} refers to an unpublished report generalising
the intersection matrix to arbitrary Eulerian digraphs, where the determinant counts
arborescences rather than circuits; we could not locate it. Lauri's own attempt at
extensions \cite{Lauri} is described by its reviewer as inconclusive. For rotor-routing
the picture is different: T\'othm\'er\'esz \cite{Toth} constructs a canonical simply
transitive action of the sandpile group on compatible Euler tours of any Eulerian ribbon
digraph, whatever the degrees, so the torsor of \cite{Bio2} generalises even though, as
we show, the matrix presentation does not.

\section{Setting}\label{sec:setting}

$D_w$ is the transition digraph of the circular word $w$: vertices the letters, one arc
$e\to f$ per occurrence of $e$ immediately followed by $f$ around the circle; loops are
deleted (they do not affect the Laplacian). It is balanced with $d^{\pm}(e)=r_e$. Its
sandpile group is $\K(D_w)=\coker(\Delta_{\hat s}^{\mathsf T})$ with $\Delta=D^{+}-T$,
$T_{ef}$ the number of transitions $e\to f$; the isomorphism class does not depend on the
sink $s$ \cite[Lem.~4.12]{HLMPPW}. Groups are compared by elementary divisors. The
\emph{pairwise pattern} of two letters $e,f$ is the circular binary word obtained by
deleting all other letters from $w$, taken up to rotation.

\section{Two families}\label{sec:two}

\begin{theorem}[Two families]\label{thm:two}
Let $w$ be a circular word over $\{A,B\}$ with $r$ copies of $A$ and $s$ of $B$, and let
$c$ be the number of transitions $A\to B$ (equivalently the number of maximal blocks of
$A$'s). Then $1\le c\le\min(r,s)$, every such $c$ occurs, and
\[
\K(D_w)\;\cong\;\Z/c .
\]
\end{theorem}

\begin{proof}
$D_w$ has two vertices; the number of arcs $A\to B$ equals the number of arcs $B\to A$
(each block of $A$'s is entered once and left once), namely $c$, and the remaining arcs
are loops. Reduced at $B$, the Laplacian is the $1\times1$ matrix $(c)$. The word
$A^{r-c+1}B^{s-c+1}(AB)^{c-1}$ realises $c$, and $c\le\min(r,s)$ since each block
contains at least one letter.
\end{proof}

For $r=s=2$ this is the interleaving theorem \cite[Thm.~6]{Bio1}: $c=1$ for
$A\,A\,B\,B$ and $A\,B\,B\,A$ (disjoint or nested, no arrangement ambiguity) and $c=2$
for $A\,B\,A\,B$. For two families the arrangement ambiguity is thus a single cyclic
factor whose order is a crossing number, and no sign is needed. Check~(A) verifies the
theorem over all $119$ circular words with $r,s\le5$.

\section{Beyond two families nothing pairwise survives}\label{sec:neg}

\begin{proposition}[Pairwise patterns do not determine the group]\label{prop:pair}
The circular words
\[
w_1=A\,A\,B\,A\,C\,C\,B\,B\,C,\qquad w_2=A\,A\,B\,A\,B\,C\,B\,C\,C,
\]
each letter three times, have the same pairwise pattern for every pair of letters, and
\[
\K(D_{w_1})=\Z/3,\qquad \K(D_{w_2})=\Z/4 .
\]
Consequently no square matrix whose $(e,f)$ entry is a function of the pairwise pattern
of $e$ and $f$ can present $\K(D_w)$ for three-copy repeats, even for $m=3$.
\end{proposition}

\begin{proof}
Restricting to $\{A,B\}$, $\{A,C\}$ and $\{B,C\}$ gives, for both words, the cyclic
patterns $001101$, $000111$ and $001101$ (with the first letter of the pair written as
$1$). In $w_1$ all six off-diagonal transition counts are $1$, so the Laplacian reduced
at $C$ is $\bigl(\begin{smallmatrix}2&-1\\-1&2\end{smallmatrix}\bigr)$, of determinant
$3$ and content $1$: $\Z/3$. In $w_2$ the transitions are $A\to B$ twice, $A\to C$
never, $B\to A$ once, $B\to C$ twice, giving
$\bigl(\begin{smallmatrix}2&-2\\-1&3\end{smallmatrix}\bigr)$, of determinant $4$ and
content $1$: $\Z/4$.
\end{proof}

\begin{remark}[Why this is sharper than the two-copy failure]\label{rem:sharper}
For two-copy repeats the pairwise pattern of $e,f$ is a single bit (interlaced or not),
and the group is likewise not a function of these bits alone --- but the first failure
is at $m=6$ \cite[Rem.~6.7]{MMN}, and it is repaired by a sign: reading from a base
point, record whether the occurrences run $e\,f\,e\,f$ or $f\,e\,f\,e$. The sign depends
on the base point, and the cokernel does not, because changing base point acts by
pivots \cite[Lem.~4.8]{MMN}. For three or more copies the analogous base-pointed pairwise
data --- the linear restriction of $w$ to each pair --- determine the whole word for
every $m$ (they fix the relative order of any two occurrences), so ``the group is a
function of them'' is vacuous, and the only genuine pairwise invariants are the rotation
classes of Proposition~\ref{prop:pair}, which fail already at $m=3$. There is no room
between the two for a sign.
\end{remark}

\begin{proposition}[No reduction to two copies]\label{prop:detach}
Let $G$ be a balanced digraph with a vertex $v$ of in- and out-degree $3$, and let $G'$
be obtained by detaching $v$ into $v'$ (two of the in-arcs, one of the out-arcs) and
$v''$ (the remaining in-arc, the other two out-arcs) joined by a new arc $v'\to v''$, so
that $G'$ is $2$-in $2$-out at both. Then $\K(G)$ is in general neither a quotient nor a
subgroup of $\K(G')$: there is a $G$ with $\K(G)=\Z/4\oplus\Z/5$ whose nine detachments
have groups $\Z/23$, $\Z/25$, $\Z/31$, $\Z/32$, $\Z/2\oplus\Z/13$, $\Z/3\oplus\Z/8$,
$\Z/2\oplus\Z/3\oplus\Z/4$ and $\Z/2\oplus\Z/3\oplus\Z/5$.
\end{proposition}

Over $150$ random balanced digraphs and all nine detachments each, $\K(G)$ is a quotient
of $\K(G')$ in $140$ of $1350$ cases (check~(C)). Euler circuits of $G'$ are those of $G$
compatible with the chosen transition at $v$, a subset of no fixed proportion, and the
group follows the count. So the arrangement group of a multi-copy skeleton cannot be
computed by splitting it into a two-copy one and applying \cite{MMN}.

\begin{remark}[What does determine it]
Trivially, the transition matrix $T$: the arrangement group is the sandpile group of the
transition digraph, which is what \cite{Bio3} says. The content of the two-copy theory is
that $T$ can be traded for interlacement --- a different, topological parametrisation of
the same group. Propositions~\ref{prop:pair} and~\ref{prop:detach} say that for three or
more copies no such trade is available through pairwise data or through two-copy
surrogates. A Gram-matrix route (the form $x^{\mathsf T}(I-S)y$ of the position cycle
restricted to consecutive-occurrence differences) was also tried and is not the
Merino--Moffatt--Noble matrix even for two copies; it is recorded in the ancillary
README so that it is not retried.
\end{remark}

\section{Census}\label{sec:census}

Table~\ref{tab:census} lists, for each multiplicity vector, the number of circular words
(one per rotation class), the number of distinct transition digraphs up to relabelling,
the number of distinct groups, the largest order, and the non-cyclic groups that occur.
The last column counts the classes of pairwise patterns and, in parentheses, how many of
them carry more than one group. For two families the groups are exactly $\Z/c$,
$c\le\min(r,s)$, as Theorem~\ref{thm:two} predicts, and every pairwise class is pure.
From three families on the arrangement group is far richer than for two-copy repeats
--- three two-copy repeats give only $0,\Z/2,\Z/3,(\Z/2)^{2}$ \cite[Prop.~8]{Bio1},
three three-copy repeats give nine groups up to order $9$ --- and most pairwise classes
are mixed.

\begin{table}[ht]
\centering\footnotesize
\begin{tabular}{lrrrrlr}
\toprule
mult. & words & $D_w$ & groups & $\max|\K|$ & non-cyclic groups & pairwise (mixed)\\
\midrule
$(3,3)$       &     4 &  3 &  3 &  3 & --- & 3 (0)\\
$(4,4)$       &    10 &  4 &  4 &  4 & --- & 7 (0)\\
$(3,3,3)$     &   188 & 12 &  9 &  9 & $(\Z/2)^{2},(\Z/3)^{2}$ & 15 (6)\\
$(3,3,2)$     &    70 & 15 &  7 &  6 & $(\Z/2)^{2}$ & 10 (4)\\
$(2,2,3)$     &    30 & 10 &  5 &  4 & $(\Z/2)^{2}$ & 6 (2)\\
$(4,4,4)$     &  2896 & 26 & 17 & 16 & $(\Z/2)^{2},(\Z/2)^{2}{\oplus}\Z/3,(\Z/3)^{2},(\Z/4)^{2},\Z/2{\oplus}\Z/4$ & 174 (86)\\
$(3,3,3,3)$   & 30804 & 96 & 28 & 27 & those, and $(\Z/2)^{3},(\Z/3)^{3},\Z/2{\oplus}(\Z/3)^{2}$ & 180 (132)\\
\bottomrule
\end{tabular}
\caption{Arrangement groups of multi-copy repeats, exhaustive over circular words
(check~(B)). $D_w$: distinct transition digraphs. All entries exact.}
\label{tab:census}
\end{table}

\section{$\varphi$X174}\label{sec:phix}

Table~\ref{tab:phix} gives the out-degree profile of the bundle-contracted skeleton of
$\varphi$X174 (NC\_001422.1) at $k=12,\dots,7$. At $k=11$ and $k=10$ every skeleton
vertex has in- and out-degree $2$, and the signed interlace matrix of an Euler circuit
found by Hierholzer's algorithm gives $\coker(A+I)=\Z/3\oplus\Z/11$ and
$\Z/4\oplus\Z/1076338579$, the groups computed in \cite{Bio3} from the Laplacian
(check~(D)) --- the theorem of \cite{MMN} on a real genome. From $k=9$ down the skeleton
has vertices of out-degree $3$ ($5$ of them at $k=9$, $42$ at $k=8$, $147$ at $k=7$),
$4$, $5$ and $6$: the two-copy presentation stops applying at the same $k$ at which the
skeleton grows from $37$ to $110$ vertices and the group acquires its large prime factors.

\begin{table}[ht]
\centering\small
\begin{tabular}{rrrl}
\toprule
$k$ & skeleton & arcs & out-degree multiset\\
\midrule
12 &   1 &    0 & (single vertex)\\
11 &   8 &   16 & $2^{8}$\\
10 &  37 &   74 & $2^{37}$\\
 9 & 110 &  227 & $2^{104}\,3^{5}\,4^{1}$\\
 8 & 312 &  679 & $2^{264}\,3^{42}\,4^{5}\,5^{1}$\\
 7 & 826 & 1945 & $2^{619}\,3^{147}\,4^{41}\,5^{12}\,6^{7}$\\
\bottomrule
\end{tabular}
\caption{$\varphi$X174: size and out-degree profile of the skeleton ($d^{\,n}$ = $n$
vertices of out-degree $d$). At $k=11,10$ the skeleton is $2$-in $2$-out and
$\coker(A+I)$ equals $\K(\Sk)$; below, multi-copy vertices appear. Exact.}
\label{tab:phix}
\end{table}

\section{Verification}\label{sec:battery}

The script \texttt{multicopy.py} runs the four checks of Table~\ref{tab:battery} in
about forty seconds, all passing; its output is archived as
\texttt{multicopy\_output.txt}. All arithmetic is exact; groups are compared by
elementary divisors.

\begin{table}[ht]
\centering\small
\begin{tabular}{clc}
\toprule
key & check & mode\\
\midrule
(A) & Theorem~\ref{thm:two} on all $119$ two-letter circular words, $r,s\le5$ & exact\\
(B) & Table~\ref{tab:census}; the witness of Proposition~\ref{prop:pair} & exact\\
(C) & Proposition~\ref{prop:detach}: $150$ digraphs, $1350$ detachments; the witness & exact\\
(D) & Table~\ref{tab:phix}; $\coker(A+I)=\K(\Sk)$ at $k=11,10$ & exact\\
\bottomrule
\end{tabular}
\caption{The verification battery, $4/4$.}
\label{tab:battery}
\end{table}

\section{Scope}\label{sec:scope}

\emph{Proved.} Theorem~\ref{thm:two}; Propositions~\ref{prop:pair}
and~\ref{prop:detach} by explicit witness.

\emph{Quoted.} The splitting theorem \cite{Bio3}; the two-copy presentation \cite{MMN};
the determinant formula \cite{MacrisPule,Jonsson} and interlace polynomial \cite{ABS};
sink-independence \cite{HLMPPW}; the rotor torsor for Eulerian ribbon digraphs
\cite{Toth}.

\emph{Open.} (a) Whether any presentation of $\K(D_w)$ for $r\ge3$ exists whose entries
are ``local'' in a weaker sense than pairwise --- e.g.\ functions of the patterns of
triples --- and whether Jonsson's unpublished generalised intersection matrix, which
counts arborescences, has a cokernel equal to $\K$. Theorem~\ref{thm:two} shows that at
least for two families the group is a crossing number, so a multi-chord crossing calculus
is not ruled out; Proposition~\ref{prop:pair} shows it cannot be pairwise. (b) The
census suggests questions we have not pursued: which groups occur for $(r,r,r)$ as $r$
grows, and whether the largest order is always attained by an elementary abelian group.
(c) The $\varphi$X174 skeleton at $k=8$ ($312$ vertices) is beyond the dense exact Smith
form used here; a modular Smith form would extend Table~\ref{tab:phix} of \cite{Bio3}.

\section*{Provenance}

Unrefereed preprint. Theorem~\ref{thm:two} is elementary and, as far as we can find,
unstated in this form; the propositions are negative results established by explicit
computation and are the reason this note exists: the natural extension of \cite{MMN} to
multi-copy repeats was the obvious next step after \cite{Bio3}, and it does not exist in
any pairwise form. We searched for an interlacement or determinant formula for Euler
circuits of $r$-in $r$-out digraphs with $r\ge3$ and found none; a search cannot
establish a negative.

\begin{thebibliography}{99}

\bibitem{ABS} R.~Arratia, B.~Bollob\'as and G.~B.~Sorkin,
\emph{The interlace polynomial of a graph},
J.\ Combin.\ Theory Ser.~B \textbf{92} (2004), 199--233.

\bibitem{Bio1} R.~Chen,
\emph{Assembly ambiguity is not a function of branch statistics: the sandpile group of a
de Bruijn graph}, companion preprint (Bio~1, version~2), 2026.

\bibitem{Bio2} R.~Chen,
\emph{The sandpile group as a coordinate system on genome reconstructions}, companion
preprint (Bio~2), 2026.

\bibitem{Bio3} R.~Chen,
\emph{The sandpile group of a de Bruijn graph splits along its repeats}, companion
preprint (Bio~3), 2026.

\bibitem{HLMPPW} A.~E.~Holroyd, L.~Levine, K.~M\'esz\'aros, Y.~Peres, J.~Propp and
D.~B.~Wilson, \emph{Chip-firing and rotor-routing on directed graphs}, in: In and Out of
Equilibrium~2, Progr.\ Probab.\ \textbf{60}, Birkh\"auser, 2008, 331--364.

\bibitem{Jonsson} J.~Jonsson,
\emph{On the number of Euler trails in directed graphs},
Math.\ Scand.\ \textbf{90} (2002), 191--214.

\bibitem{Lauri} J.~Lauri,
\emph{On a determinant formula for enumerating Euler trails in a class of digraphs},
Rend.\ Sem.\ Mat.\ Messina Ser.~II \textbf{5} (1999), 75--90.

\bibitem{MacrisPule} N.~Macris and J.~V.~Pul\'e,
\emph{An alternative formula for the number of Euler trails for a class of digraphs},
Discrete Math.\ \textbf{154} (1996), 301--305.

\bibitem{MMN} C.~Merino, I.~Moffatt and S.~Noble,
\emph{The critical group of a combinatorial map},
Combinatorial Theory \textbf{5} (2025), no.~3; arXiv:2308.13342.

\bibitem{Toth} L.~T\'othm\'er\'esz,
\emph{On canonical sandpile actions of embedded graphs},
Advances in Combinatorics (2026), no.~6; arXiv:2504.05275.

\end{thebibliography}

\end{document}
